\section{Objetivos}\label{objetivos}

Este apartado constituye el puente entre el problema identificado en el capítulo anterior y la contribución que se propone. El esquema de reabastecimiento periódico y de rutas fijas que utiliza Cafetalino provoca visitas innecesarias a equipos que aún disponen de insumos y, al mismo tiempo, episodios de desabastecimiento que interrumpen el servicio, de modo que la propuesta debe traducirse en un efecto medible sobre esa operación. En consecuencia, a continuación se presentan tanto el objetivo general, como los objetivos específicos que contribuyen a su logro, planteados para el presente proyecto de innovación.

\subsection{Objetivo General}\label{objetivo-general}

Mejorar la gestión del reabastecimiento de la red de máquinas dispensadoras de Cafetalino, reduciendo las visitas innecesarias y los episodios de desabastecimiento, mediante un sistema basado en Inteligencia Artificial que prediga la demanda y permita la planificación dinámica y la optimización de las rutas logísticas, con el fin de contribuir a la continuidad del servicio, la reducción de ineficiencias operativas, mejorar la experiencia del cliente y minimizar los costos logísticos.

\subsection{Objetivos Específicos}\label{objetivos-especuxedficos}

Para el logro del objetivo general, se formularon los siguientes objetivos específicos:

\begin{itemize}
\item
  Caracterizar los patrones históricos de consumo y reabastecimiento de la red de máquinas dispensadoras de Cafetalino a partir de la información operativa disponible entre 2018 y 2026.
\item
  Desarrollar un módulo para generar estimaciones de los requerimientos futuros de reposición de insumos mediante la aplicación de técnicas de aprendizaje automático basadas en datos históricos de operación, sin la instalación de hardware adicional.
\item
  Desarrollar un módulo para optimizar la planificación de recorridos de abastecimiento mediante la utilización de rutas dinámicas, construidas a partir de la demanda proyectada para cada máquina dispensadora.
\item
  Desarrollar una interfaz web para el personal operativo de Cafetalino, que se constituya en una herramienta de apoyo a la toma de decisiones, y permita visualizar pronósticos de la demanda y las rutas de abastecimiento recomendadas.
\item
  Evaluar el desempeño de la solución propuesta mediante métricas de error de predicción y de reducción de recorridos frente al esquema de rutas fijas actual.
\end{itemize}
